\documentclass[12pt]{article}
\usepackage[T1]{fontenc}
\usepackage[utf8]{inputenc}
\usepackage{lmodern}
\usepackage{textcomp}
\usepackage{enumitem}
\usepackage{geometry}
\geometry{margin=1in}
\begin{document}
\begin{center}
Answer Key - Sociology - Graduate Level Assessment 19 \
\small (Answers are ordered exactly as in the question paper.)
\end{center}

\section*{Multiple Choice}
\begin{enumerate}[label=\arabic*.]
\item \textbf{Answer:} C
\\ \textit{Options:} A) there are no clearly defined projects into which the money can be directed; B) the United Nations has refused to call on rich countries to provide it; C) debt repayments with interest can be greater than the amount of money received; D) debt repayments with interest can be greater than the amount of money received
\\ \textit{Note:} The correct answer highlights that the financial burden of debt repayments, often coupled with high interest rates, can negate the benefits of economic aid received, leaving developing countries in a cycle of dependency and hindering their modernization efforts.
\item \textbf{Answer:} B
\\ \textit{Options:} A) pollies; B) mollies; C) dollies; D) lollies
\\ \textit{Note:} The term "mollies" refers specifically to a deviant subculture of homosexual men in seventeenth and eighteenth century London, characterized by their distinct social practices and gatherings, thus making it the correct answer to the question.
\item \textbf{Answer:} A
\\ \textit{Options:} A) work can be subcontracted out to independent suppliers and retailers; B) business transactions occur only through electronic communication; C) the Japanese model is applied, through lateral networks of flexible roles; D) activities are redistributed equally between men and women
\\ \textit{Note:} The correct answer reflects the key concept of network organization, which emphasizes flexibility and efficiency by allowing businesses to outsource specific tasks to independent suppliers and retailers, thereby enhancing collaboration and resource optimization.
\item \textbf{Answer:} C
\\ \textit{Options:} A) effervescent ceremonies that create a feeling of belonging; B) images of gods or totems that are widely recognized; C) shared ideas and moral values, often symbolized by an object or figurehead; D) ideological tools used to obscure class divisions
\\ \textit{Note:} The correct answer highlights that 'collective representations' in Durkheim's work encapsulate the shared beliefs and moral values of a society, which are often embodied and communicated through symbols or figures that unify and guide collective social behavior.
\item \textbf{Answer:} A
\\ \textit{Options:} A) self-governing trusts competing for purchasing contracts from health authorities; B) state-controlled providers, dependent on funding from the central government; C) increasingly detached from health authorities and providers of private health care; D) less inclined to run themselves efficiently, as demand for health care was falling
\\ \textit{Note:} The 1990 reform of the National Health Service established hospitals as self-governing trusts, allowing them to operate independently and compete for purchasing contracts from health authorities, which was a significant shift towards market-oriented principles in healthcare management.
\end{enumerate}

\section*{True / False}
\begin{enumerate}[label=\arabic*.]
\setcounter{enumi}{5}
\item \textbf{Answer:} True
\\ \textit{Options:} A) True; B) False
\\ \textit{Note:} Pre-testing a questionnaire allows researchers to identify and resolve ambiguities in question wording and routing issues, ensuring that the data collected is valid and reliable.
\item \textbf{Answer:} False
\\ \textit{Options:} A) True; B) False
\\ \textit{Note:} The lack of transport systems in the nineteenth century actually hindered urbanization, as individuals in rural areas relied on agriculture and local markets, making city living less practical due to limited access and connectivity.
\item \textbf{Answer:} True
\\ \textit{Options:} A) True; B) False
\\ \textit{Note:} Modernization theory asserts that societies evolve through a linear process that transitions from pre-industrial, traditional structures to advanced, industrialized forms, reflecting a progression of social, economic, and technological development.
\item \textbf{Answer:} False
\\ \textit{Options:} A) True; B) False
\\ \textit{Note:} The statement is false because the ruling elite can include individuals from diverse class backgrounds and experiences, reflecting a variety of perspectives and interests in political representation.
\item \textbf{Answer:} False
\\ \textit{Options:} A) True; B) False
\\ \textit{Note:} The post-war welfare state of 1948 primarily aimed to provide a social safety net through social security and unemployment benefits rather than guaranteeing free health care and education for all citizens, as these services became more widely available in subsequent decades.
\end{enumerate}

\section*{Long Form Response}
\begin{enumerate}[label=\arabic*.]
\setcounter{enumi}{10}
\item \textbf{Answer:} Socialization is the lifelong process through which individuals acquire the norms, values, behaviors, and social skills necessary for functioning within their society. It begins in early childhood and continues throughout life, significantly shaping individual identity by integrating personal experiences with cultural contexts. Through agents of socialization such as family, peers, education, and media, individuals learn to navigate social structures and develop a sense of self that aligns with societal expectations. This process not only fosters individual identity but also facilitates membership in society by promoting social cohesion and continuity, as individuals internalize shared beliefs and practices. Ultimately, socialization is crucial for both personal development and the perpetuation of social order.
\\ \textit{Note:} correct because it effectively outlines the lifelong nature of socialization, identifies key agents that influence identity formation, and emphasizes the integration of personal experiences with cultural contexts, all of which are essential for understanding how individuals function within society.
\item \textbf{Answer:} The functionalist perspective in sociology conceptualizes society as a complex system of interdependent parts, where each component plays a vital role in maintaining the overall stability and functionality of the social order. According to this perspective, social institutions such as family, education, and religion work together to fulfill essential functions, ensuring cohesion and continuity within the society. This interdependence implies that any significant change in one part can lead to adjustments in others, highlighting the delicate balance required for social stability. However, functionalism also recognizes that social change is inevitable and can occur through gradual evolution or significant disruption, often prompted by shifts in values, technology, or demographic trends. Understanding society through a functionalist lens thus enables a nuanced analysis of both the mechanisms that uphold stability and the processes that facilitate change.
\\ \textit{Note:} correct because it accurately describes the functionalist perspective's view of society as a system of interdependent parts that work together to maintain stability, emphasizing the significance of social institutions in fulfilling essential functions and the potential impact of change within this system.
\end{enumerate}

\vfill
\noindent\textit{End of Answer Key}
\end{document}